\begin{tablalafrox}{Diccionario de la cuenta 2.3. Fuente: (Bases Administrativas, art. 17, pp. 12–13, y Formulario T-14, p. 63); (Bases Técnicas del caso, secciones 17.2–17.6, pp. 32–34).}{tab:T14-dic-2-3}
  {L{1.2cm} Y Y Y C{0.9cm} L{2.3cm}}
  {\cab{Código} & \cab{Paquete de trabajo} & \cab{Entregable} & \cab{Criterio de aceptación} & \cab{Resp.} & \cab{Período}}
  2.3.1 & Plan de seguridad de la información & Plan de seguridad. & Está aprobado en el H2, y cada requisito de seguridad del T-12 tiene su control. & SEG & Ambas · Mes 2 (H2). \\
  2.3.2 & Modelado de amenazas por componente y por integración externa, actualizado ante cada cambio & Modelo de amenazas con sus controles. & Cada amenaza alta tiene un control probado o un riesgo aceptado por escrito por el CLIENTE. & SEG & Ambas · Meses 3 a 10, y ante cada cambio de arquitectura. \\
  2.3.3 & Plan de respuesta a incidentes de seguridad, con aviso al CLIENTE en dos horas & Plan de respuesta a incidentes de seguridad. & El plan define la clasificación, el escalamiento, los plazos y los responsables, y se ensaya una vez antes del H5. & SEG & Ambas · Mes 3. \\
  2.3.4 & Matriz de clasificación de la información y matriz de controles & Las dos matrices. & Cada tipo de dato tiene su nivel, y cada nivel tiene sus controles implementados. & SEG & Ambas · Mes 2 (con el H2). \\
\end{tablalafrox}
